\clearpage
\section{The checked Lean argument: objects and dependency map}\label{app:lean-map}

This appendix presents the formal argument inside the paper. It explains how
the source objects are represented and prints the complete central declaration
bodies used in the distance obstruction, together with the applications needed
to recover the mathematical theorems. Appendix~\ref{app:lean-links} then gives
current and immutable links for every admitted source file. These links
supplement the displayed proofs; they are also the route to routine coordinate
lemmas and the large finite row certificates not printed in full.

All listings below are generated from the checked sources at commit
\texttt{98eb91cd9049}. They retain complete declarations, with Unicode symbols
typeset mathematically and long lines wrapped. They are extracts from their
original namespaces and imported contexts, rather than separate standalone Lean
files. Unless stated otherwise, cyclic modules use an implicit natural number
\texttt{L} and the instance \texttt{[NeZero L]}. This excludes length zero, not
even or composite lengths. The generic cover modules instead use finite abelian
groups and the instances stated with their theorem signatures. The archive's
\nolinkurl{appendix-listings.json} binds each extract to its source hash and exact
line ranges. No mathematical proof source was changed to produce these listings.

\subsection{How the branches fit together}
The lower and upper bounds are proved independently and combined only after
their actual source-coordinate translations. The following diagram summarizes
the dependency structure; the graph and transpose branches establish the
additional clauses of Theorem~\ref{thm:main}.
\[
\begin{array}{ccc}
\text{Kernel and augmentation}&&\text{Base certificate and tree bound}\\
\downarrow&&\downarrow\\
\text{Full-span exclusion; weight }\leq11&&\text{Cover projection; weight }\geq12\\
\searrow&&\swarrow\\[-2pt]
\multicolumn{3}{c}{\text{Actual source words, spans and minima}}\\
\multicolumn{3}{c}{\downarrow}\\
\multicolumn{3}{c}{d_Q\leq11<12\leq d_C}
\end{array}
\]
\begin{center}\small
\begin{tabular}{@{}P{.30\textwidth}cP{.57\textwidth}@{}}
\toprule Branch&&Checked dependency chain\\\midrule
Quantum upper bound&$\longrightarrow$&
Coefficient kernel $\to$ augmentation separator $\to$ full-span nonmembership
$\to$ physical support at most eleven.\\[4pt]
Classical lower bound&$\longrightarrow$&
Base row reconstruction and finite weight certificate; Tanner-tree stopping bound
$\to$ parity projection $\to$ every nonzero upstairs word has weight at least twelve.\\[4pt]
Source distances&$\longrightarrow$&
Actual matrix actions and full spans $\to$ both CSS sectors and all-word weight
predicates $\to$ equality of the formal minima $\to$ $d_Q<d_C$.\\[4pt]
Graph and transpose properties&$\longrightarrow$&
Voltage exclusions plus explicit cycle edges; cyclic reachability;
transpose period identity $\to$ constant fibers $\to$ exact transpose weight.\\
\bottomrule
\end{tabular}
\end{center}
The trust boundary is the ordinary Lean kernel and the recorded standard axioms
\texttt{propext}, \texttt{Classical.choice} and \texttt{Quot.sound}. The finite
classical certificate is checked by Lean; external C++ and SAT exclusions are
absent from these branches. The separately audited generic follow-up lemmas are
not used to claim the unproved truncated-distance condition.

\subsection{Actual matrices and physical coordinates}
\nolinkurl{GeneralLP} gives the coefficient matrices for an arbitrary monomial
$3\times4$ exponent array. Its \texttt{Qubit} type is the disjoint union of the
$4\times4$ and $3\times3$ product blocks, each with one cyclic coordinate.
The \texttt{P} coefficient uses exactly $t=s+e$. Thus the definitions implement
\eqref{eq:source}, including the star's negative exponent, without a second lift.
The syndrome and two stabilizer actions displayed here expose every sign.
\leanlisting{source-objects}

\subsection{The complete row span, not just the generators}
The following equivalence quantifies over an arbitrary binary coefficient for
every lifted check. Its forward direction evaluates a general row combination
in each physical block; its reverse direction reconstructs that combination.
The \texttt{inSpanZ\_iff} theorem in the same file proves the corresponding
other-sector statement by transposing both square blocks. This is the bridge
needed before an augmentation obstruction can certify a logical class.
\leanlisting{source-span}

\subsection{The all-word distance predicate}
Here \texttt{Kernel} and \texttt{InSpan} quantify over actual source binary
coordinates, while \texttt{QuantumHasWeight} includes both logical sectors.
These are the predicates behind the minima, rather than a restricted search
through displayed words. Existence and the minimum equivalences are used in
Appendix~\ref{app:lean-assembly}.
\leanlisting{source-minimum}

\clearpage
\section{The logical obstruction: kernel, full-span separation and weight}\label{app:lean-obstruction}
This branch proves Theorem~\ref{thm:template}. The coefficient word is defined
for arbitrary parameters in \texttt{ZMod L}. In particular, monomial terms may
coincide, and the weight argument must allow characteristic-two cancellation.

\subsection{The word and its full coefficient kernel proof}
The matrix literal below is the coefficient version of \eqref{eq:UV}.
\texttt{delta a s} is one exactly when $s=a$. The proof checks all twelve
check types with arbitrary fiber coordinate; it does not sample lift sizes.
\leanlisting{parametric-word}

\subsection{Why the separator annihilates every stabilizer}
The augmentation is the two fiber sums of \eqref{eq:ell}. Translation permutes
a finite cyclic fiber, so both terms of a stabilizer give the same sum and
cancel. The second theorem extracts a left-block stabilizer representation
from membership in the entire physical row span. Together these statements
make the separator apply to arbitrary source stabilizer combinations.
\leanlisting{separator-span}

\subsection{Nontriviality of the actual source logical}
The source check action transports the coefficient identity to the matrix
kernel. If the word belonged to the stabilizer span, its left block would
have augmentation zero, contradicting the computed value one. The following
is the complete proof of that nonmembership, with no finite span search.
\leanlisting{logical-separation}

\subsection{Physical Hamming weight, including cancelling parameters}
The support list consists of eleven physical qubit coordinates. The first
proof shows the word vanishes outside that list; the second compares finite
support cardinalities. It deliberately uses list length as an upper bound,
so duplicate coordinates and cancellation cannot invalidate the theorem.
\leanlisting{physical-weight}

\subsection{The source-matched universal conclusion}
This theorem assembles actual CSS matrix commutation, full logical
nontriviality and physical coefficient weight. It is the same universal
parameter statement used in Theorem~\ref{thm:template}.
\leanlisting{source-obstruction}

\clearpage
\section{The classical lower bound: projection and the finite base}\label{app:lean-classical}
This branch separates ordinary binary parity projection from the image of
occupied support. The distinction is necessary when all projected parities
cancel, and avoids an odd-cover assumption.

\subsection{Cancellation gives two occupied coordinates per fiber}
The push-forward is a sum in \texttt{ZMod 2}. If it is zero, a nonempty
fiber support cannot have cardinality one. Summing the resulting two-per-fiber
bound gives twice the size of the occupied image. The full proof below
works for arbitrary finite source and target types, not just cyclic groups.
\leanlisting{cancelled-fibers}

\subsection{The occupied image is a nonempty stopping set}
The projected support theorem uses the actual monomial incidence. At a
touched lifted check, binary parity gives a second selected neighbor.
Its base column type differs from the first, so the neighbors remain
distinct after projection. The helper \texttt{stopping\_of\_partners}
turns these partners into the finite-set stopping predicate.
\leanlisting{stopping-projection}

\subsection{The complete cover-bound theorem}
Both alternatives now appear in a single full theorem and proof. When
the parity projection is nonzero, its base code distance supplies $d$;
when it is zero, the occupied image supplies $2s$. The printed hypotheses
make clear which base facts are required.
\leanlisting{cover-bound}

\subsection{What the eleven-fold certificate actually checks}
\nolinkurl{FamilyEnum11.encode} is an explicit XOR of thirteen $44$-bit
vectors. \nolinkurl{FamilySource11} proves
\texttt{kernel\_reconstruct}: every actual source-kernel vector equals
that encoding at its thirteen extracted free coordinates. It is an
identity derived from the source rows, not an assumed generator-span
description. \nolinkurl{FamilyEnum11.goodBlock} checks 128 consecutive
inputs, and \nolinkurl{FamilyDistance11.all_blocks} checks all 64 blocks.
Their complete source files, including the bit constants and row identities,
are linked in Appendix~\ref{app:lean-links}.

The following proofs turn the checked blocks into a theorem for an arbitrary
thirteen-bit input and then for every source-kernel vector. This displays
the certificate's soundness step rather than an external enumerator result.
\leanlisting{finite-certificate}

The next bridge packs actual coefficient functions into the bit-vector
certificate and identifies population count with physical Hamming weight.
The final theorem converts each source check equation into a certified row
dot product before applying the lower bound.
\leanlisting{finite-weight-bridge}

\subsection{The structural stopping bound}
The classical Tanner-tree argument is credited to Kelley et al.~\cite{Kelley}
in the main text. Here is its complete checked degree-three, no-four/no-six
implementation. It constructs six distinct selected variables explicitly;
the no-cycle hypotheses rule out each potential collision.
\leanlisting{tree-stopping}

\subsection{Application to every positive multiple of eleven}
Reduction modulo eleven supplies the additive map. The two premises of
the generic theorem are discharged by the actual base kernel certificate
and the independently proved structural stopping bound. This is the full
family application, including the exponent-projection equality.
\leanlisting{cover-application}

\clearpage
\section{Upper bounds, transpose distance and actual graph witnesses}\label{app:lean-geometry}

\subsection{A nonzero classical word of weight at most twenty}
The coefficient definition below is \eqref{eq:twenty}. Kernel membership
is proved for every positive $L$; nonzero membership uses reduction modulo
eleven. The final step applies the physical support bound to the actual
classical minimum. The helper \texttt{weight\_le} and its support list
are in the same linked source file.
\leanlisting{classical-upper}

\subsection{The transpose period identity}
The transpose boundary has the plus exponent $t+E_{ij}$. The following
proof writes down the precise combination of shifted check equations in
\eqref{eq:bezout}. The identity forces the third component to have period
one without semisimplicity, oddness or a rank-preservation hypothesis.
\leanlisting{transpose-period}

The intervening lemmas in \nolinkurl{TransposeConstant.lean} propagate
period one to constancy of all three fibers. The full equivalence below
includes the converse and the even-sum condition.
\leanlisting{transpose-constant}

\subsection{Exact weight of every nonzero transpose word}
The finite binary-triple fact \texttt{three\_even\_count} says that a
nonzero even triple has exactly two ones. The proof below then counts
the Cartesian product of those two occupied component types with all
$L$ fibers. The final conjunction includes existence, so the distance
statement is not vacuous.
\leanlisting{transpose-distance}

\subsection{Girth exactly eight}
The finite voltage tests and their reduction proofs exclude four- and
six-cycles. To complete exact girth, one also needs a simple eight-cycle
in the actual incidence. The following declarations give distinct
vertices and all eight edges separately. The main text explains their
assembly into the displayed cycle.
\leanlisting{girth-eight}

\subsection{Connectedness through actual Tanner walks}
The normalized zero row connects column types within a fiber; the unit
exponent supplies the following shift walk. The final proof combines
variable-to-variable walks and the zero-column incidences to cover all
variable/check endpoints. Both statements hold for every positive $L$.
\leanlisting{connected}

\clearpage
\section{Parameter transfer and the final source-distance statement}\label{app:lean-assembly}

\subsection{The arithmetic counterfamily}
This complete implication displays every premise of the general parameter
extension. An additive projection sends the parameters to $(1,5)$; the
conclusion includes the voltage exclusions, full partition exclusions,
actual classical-kernel lower bound and source logical upper witness.
The cyclic-congruence specialization is in the same source file.
\leanlisting{arithmetic-family}

\subsection{From block coordinates to every physical binary word}
The disjoint-union coordinate map is surjective, so every physical word
has two blocks. The first two equivalences transport logical-sector
membership and exact support weight. The next two proofs identify the
minima in both directions, using the corresponding weight equivalences.
This is the bridge that prevents an accidentally restricted minimum.
\leanlisting{minimum-bridge}

\subsection{The final distance drop}
All substantive ingredients now occur as checked dependencies. The
last theorem rewrites the actual source minima, applies the quantum
witness and classical lower theorem, and concludes strict separation.
The other clauses of Theorem~\ref{thm:main} are supplied by the graph,
transpose and upper-word proofs above; they are not extra assumptions
of this distance statement.
\leanlisting{final-distance}

\paragraph{Reproducing the checks.}
The fifty original modules and their full import closure were rebuilt
against the pinned cached dependencies as described in
Section~\ref{sec:formal}. The manuscript declaration audit checks every
named declaration linked here, including those added to the exposition.
The listing generator verifies that each source equals its published
Git blob, and that each printed declaration is a complete source range.
The archive contains that generator, the range/hash manifest, the
declaration audit and the portable rebuild instructions. The exact
typesetting extracts add no axioms or proof modules.

\clearpage\section{Complete current and immutable Lean URLs}\label{app:lean-links}
The following index includes every admitted proof, certificate and research import.
All immutable links use commit \texttt{98eb91cd9049e6a88887a5ba79150ad96a4a3e53}.
The declaration links in the text and preceding appendices go directly to source line ranges.
Current links may evolve; the immutable sources match the audited hashes.
The source archive and companion proof guide give the complete SHA-256 ledger and logs.
\par\smallskip\noindent\begin{minipage}{\linewidth}\small\setlength{\baselineskip}{11pt}
\textbf{Source file:} \nolinkurl{BinaryFinite.lean}\par\nopagebreak
Current: \url{https://github.com/alectionary/vibe-math/blob/master/formalization/distance-preservation/BinaryFinite.lean}\\
Immutable: \url{https://github.com/alectionary/vibe-math/blob/98eb91cd9049e6a88887a5ba79150ad96a4a3e53/formalization/distance-preservation/BinaryFinite.lean}\par\end{minipage}\par
\par\smallskip\noindent\begin{minipage}{\linewidth}\small\setlength{\baselineskip}{11pt}
\textbf{Source file:} \nolinkurl{BinaryWords.lean}\par\nopagebreak
Current: \url{https://github.com/alectionary/vibe-math/blob/master/formalization/distance-preservation/BinaryWords.lean}\\
Immutable: \url{https://github.com/alectionary/vibe-math/blob/98eb91cd9049e6a88887a5ba79150ad96a4a3e53/formalization/distance-preservation/BinaryWords.lean}\par\end{minipage}\par
\par\smallskip\noindent\begin{minipage}{\linewidth}\small\setlength{\baselineskip}{11pt}
\textbf{Source file:} \nolinkurl{CSS11.lean}\par\nopagebreak
Current: \url{https://github.com/alectionary/vibe-math/blob/master/formalization/distance-preservation/CSS11.lean}\\
Immutable: \url{https://github.com/alectionary/vibe-math/blob/98eb91cd9049e6a88887a5ba79150ad96a4a3e53/formalization/distance-preservation/CSS11.lean}\par\end{minipage}\par
\par\smallskip\noindent\begin{minipage}{\linewidth}\small\setlength{\baselineskip}{11pt}
\textbf{Source file:} \nolinkurl{ClassicalDistance11.lean}\par\nopagebreak
Current: \url{https://github.com/alectionary/vibe-math/blob/master/formalization/distance-preservation/ClassicalDistance11.lean}\\
Immutable: \url{https://github.com/alectionary/vibe-math/blob/98eb91cd9049e6a88887a5ba79150ad96a4a3e53/formalization/distance-preservation/ClassicalDistance11.lean}\par\end{minipage}\par
\par\smallskip\noindent\begin{minipage}{\linewidth}\small\setlength{\baselineskip}{11pt}
\textbf{Source file:} \nolinkurl{ClassicalEnum11.lean}\par\nopagebreak
Current: \url{https://github.com/alectionary/vibe-math/blob/master/formalization/distance-preservation/ClassicalEnum11.lean}\\
Immutable: \url{https://github.com/alectionary/vibe-math/blob/98eb91cd9049e6a88887a5ba79150ad96a4a3e53/formalization/distance-preservation/ClassicalEnum11.lean}\par\end{minipage}\par
\par\smallskip\noindent\begin{minipage}{\linewidth}\small\setlength{\baselineskip}{11pt}
\textbf{Source file:} \nolinkurl{ClassicalUpper20.lean}\par\nopagebreak
Current: \url{https://github.com/alectionary/vibe-math/blob/master/formalization/distance-preservation/ClassicalUpper20.lean}\\
Immutable: \url{https://github.com/alectionary/vibe-math/blob/98eb91cd9049e6a88887a5ba79150ad96a4a3e53/formalization/distance-preservation/ClassicalUpper20.lean}\par\end{minipage}\par
\par\smallskip\noindent\begin{minipage}{\linewidth}\small\setlength{\baselineskip}{11pt}
\textbf{Source file:} \nolinkurl{Counterexample11.lean}\par\nopagebreak
Current: \url{https://github.com/alectionary/vibe-math/blob/master/formalization/distance-preservation/Counterexample11.lean}\\
Immutable: \url{https://github.com/alectionary/vibe-math/blob/98eb91cd9049e6a88887a5ba79150ad96a4a3e53/formalization/distance-preservation/Counterexample11.lean}\par\end{minipage}\par
\par\smallskip\noindent\begin{minipage}{\linewidth}\small\setlength{\baselineskip}{11pt}
\textbf{Source file:} \nolinkurl{CoverDistance.lean}\par\nopagebreak
Current: \url{https://github.com/alectionary/vibe-math/blob/master/formalization/distance-preservation/CoverDistance.lean}\\
Immutable: \url{https://github.com/alectionary/vibe-math/blob/98eb91cd9049e6a88887a5ba79150ad96a4a3e53/formalization/distance-preservation/CoverDistance.lean}\par\end{minipage}\par
\par\smallskip\noindent\begin{minipage}{\linewidth}\small\setlength{\baselineskip}{11pt}
\textbf{Source file:} \nolinkurl{CoverWeight.lean}\par\nopagebreak
Current: \url{https://github.com/alectionary/vibe-math/blob/master/formalization/distance-preservation/CoverWeight.lean}\\
Immutable: \url{https://github.com/alectionary/vibe-math/blob/98eb91cd9049e6a88887a5ba79150ad96a4a3e53/formalization/distance-preservation/CoverWeight.lean}\par\end{minipage}\par
\par\smallskip\noindent\begin{minipage}{\linewidth}\small\setlength{\baselineskip}{11pt}
\textbf{Source file:} \nolinkurl{CyclicCover.lean}\par\nopagebreak
Current: \url{https://github.com/alectionary/vibe-math/blob/master/formalization/distance-preservation/CyclicCover.lean}\\
Immutable: \url{https://github.com/alectionary/vibe-math/blob/98eb91cd9049e6a88887a5ba79150ad96a4a3e53/formalization/distance-preservation/CyclicCover.lean}\par\end{minipage}\par
\par\smallskip\noindent\begin{minipage}{\linewidth}\small\setlength{\baselineskip}{11pt}
\textbf{Source file:} \nolinkurl{FamilyBaseDistance.lean}\par\nopagebreak
Current: \url{https://github.com/alectionary/vibe-math/blob/master/formalization/distance-preservation/FamilyBaseDistance.lean}\\
Immutable: \url{https://github.com/alectionary/vibe-math/blob/98eb91cd9049e6a88887a5ba79150ad96a4a3e53/formalization/distance-preservation/FamilyBaseDistance.lean}\par\end{minipage}\par
\par\smallskip\noindent\begin{minipage}{\linewidth}\small\setlength{\baselineskip}{11pt}
\textbf{Source file:} \nolinkurl{FamilyBaseGraph.lean}\par\nopagebreak
Current: \url{https://github.com/alectionary/vibe-math/blob/master/formalization/distance-preservation/FamilyBaseGraph.lean}\\
Immutable: \url{https://github.com/alectionary/vibe-math/blob/98eb91cd9049e6a88887a5ba79150ad96a4a3e53/formalization/distance-preservation/FamilyBaseGraph.lean}\par\end{minipage}\par
\par\smallskip\noindent\begin{minipage}{\linewidth}\small\setlength{\baselineskip}{11pt}
\textbf{Source file:} \nolinkurl{FamilyBounds.lean}\par\nopagebreak
Current: \url{https://github.com/alectionary/vibe-math/blob/master/formalization/distance-preservation/FamilyBounds.lean}\\
Immutable: \url{https://github.com/alectionary/vibe-math/blob/98eb91cd9049e6a88887a5ba79150ad96a4a3e53/formalization/distance-preservation/FamilyBounds.lean}\par\end{minipage}\par
\par\smallskip\noindent\begin{minipage}{\linewidth}\small\setlength{\baselineskip}{11pt}
\textbf{Source file:} \nolinkurl{FamilyConnected.lean}\par\nopagebreak
Current: \url{https://github.com/alectionary/vibe-math/blob/master/formalization/distance-preservation/FamilyConnected.lean}\\
Immutable: \url{https://github.com/alectionary/vibe-math/blob/98eb91cd9049e6a88887a5ba79150ad96a4a3e53/formalization/distance-preservation/FamilyConnected.lean}\par\end{minipage}\par
\par\smallskip\noindent\begin{minipage}{\linewidth}\small\setlength{\baselineskip}{11pt}
\textbf{Source file:} \nolinkurl{FamilyCoverAlternative.lean}\par\nopagebreak
Current: \url{https://github.com/alectionary/vibe-math/blob/master/formalization/distance-preservation/FamilyCoverAlternative.lean}\\
Immutable: \url{https://github.com/alectionary/vibe-math/blob/98eb91cd9049e6a88887a5ba79150ad96a4a3e53/formalization/distance-preservation/FamilyCoverAlternative.lean}\par\end{minipage}\par
\par\smallskip\noindent\begin{minipage}{\linewidth}\small\setlength{\baselineskip}{11pt}
\textbf{Source file:} \nolinkurl{FamilyDistance11.lean}\par\nopagebreak
Current: \url{https://github.com/alectionary/vibe-math/blob/master/formalization/distance-preservation/FamilyDistance11.lean}\\
Immutable: \url{https://github.com/alectionary/vibe-math/blob/98eb91cd9049e6a88887a5ba79150ad96a4a3e53/formalization/distance-preservation/FamilyDistance11.lean}\par\end{minipage}\par
\par\smallskip\noindent\begin{minipage}{\linewidth}\small\setlength{\baselineskip}{11pt}
\textbf{Source file:} \nolinkurl{FamilyDotRows11.lean}\par\nopagebreak
Current: \url{https://github.com/alectionary/vibe-math/blob/master/formalization/distance-preservation/FamilyDotRows11.lean}\\
Immutable: \url{https://github.com/alectionary/vibe-math/blob/98eb91cd9049e6a88887a5ba79150ad96a4a3e53/formalization/distance-preservation/FamilyDotRows11.lean}\par\end{minipage}\par
\par\smallskip\noindent\begin{minipage}{\linewidth}\small\setlength{\baselineskip}{11pt}
\textbf{Source file:} \nolinkurl{FamilyEnum11.lean}\par\nopagebreak
Current: \url{https://github.com/alectionary/vibe-math/blob/master/formalization/distance-preservation/FamilyEnum11.lean}\\
Immutable: \url{https://github.com/alectionary/vibe-math/blob/98eb91cd9049e6a88887a5ba79150ad96a4a3e53/formalization/distance-preservation/FamilyEnum11.lean}\par\end{minipage}\par
\par\smallskip\noindent\begin{minipage}{\linewidth}\small\setlength{\baselineskip}{11pt}
\textbf{Source file:} \nolinkurl{FamilyGirth.lean}\par\nopagebreak
Current: \url{https://github.com/alectionary/vibe-math/blob/master/formalization/distance-preservation/FamilyGirth.lean}\\
Immutable: \url{https://github.com/alectionary/vibe-math/blob/98eb91cd9049e6a88887a5ba79150ad96a4a3e53/formalization/distance-preservation/FamilyGirth.lean}\par\end{minipage}\par
\par\smallskip\noindent\begin{minipage}{\linewidth}\small\setlength{\baselineskip}{11pt}
\textbf{Source file:} \nolinkurl{FamilyNoShort.lean}\par\nopagebreak
Current: \url{https://github.com/alectionary/vibe-math/blob/master/formalization/distance-preservation/FamilyNoShort.lean}\\
Immutable: \url{https://github.com/alectionary/vibe-math/blob/98eb91cd9049e6a88887a5ba79150ad96a4a3e53/formalization/distance-preservation/FamilyNoShort.lean}\par\end{minipage}\par
\par\smallskip\noindent\begin{minipage}{\linewidth}\small\setlength{\baselineskip}{11pt}
\textbf{Source file:} \nolinkurl{FamilyRing.lean}\par\nopagebreak
Current: \url{https://github.com/alectionary/vibe-math/blob/master/formalization/distance-preservation/FamilyRing.lean}\\
Immutable: \url{https://github.com/alectionary/vibe-math/blob/98eb91cd9049e6a88887a5ba79150ad96a4a3e53/formalization/distance-preservation/FamilyRing.lean}\par\end{minipage}\par
\par\smallskip\noindent\begin{minipage}{\linewidth}\small\setlength{\baselineskip}{11pt}
\textbf{Source file:} \nolinkurl{FamilySource11.lean}\par\nopagebreak
Current: \url{https://github.com/alectionary/vibe-math/blob/master/formalization/distance-preservation/FamilySource11.lean}\\
Immutable: \url{https://github.com/alectionary/vibe-math/blob/98eb91cd9049e6a88887a5ba79150ad96a4a3e53/formalization/distance-preservation/FamilySource11.lean}\par\end{minipage}\par
\par\smallskip\noindent\begin{minipage}{\linewidth}\small\setlength{\baselineskip}{11pt}
\textbf{Source file:} \nolinkurl{FamilyStopping11.lean}\par\nopagebreak
Current: \url{https://github.com/alectionary/vibe-math/blob/master/formalization/distance-preservation/FamilyStopping11.lean}\\
Immutable: \url{https://github.com/alectionary/vibe-math/blob/98eb91cd9049e6a88887a5ba79150ad96a4a3e53/formalization/distance-preservation/FamilyStopping11.lean}\par\end{minipage}\par
\par\smallskip\noindent\begin{minipage}{\linewidth}\small\setlength{\baselineskip}{11pt}
\textbf{Source file:} \nolinkurl{FamilyTheorem.lean}\par\nopagebreak
Current: \url{https://github.com/alectionary/vibe-math/blob/master/formalization/distance-preservation/FamilyTheorem.lean}\\
Immutable: \url{https://github.com/alectionary/vibe-math/blob/98eb91cd9049e6a88887a5ba79150ad96a4a3e53/formalization/distance-preservation/FamilyTheorem.lean}\par\end{minipage}\par
\par\smallskip\noindent\begin{minipage}{\linewidth}\small\setlength{\baselineskip}{11pt}
\textbf{Source file:} \nolinkurl{FamilyTranspose.lean}\par\nopagebreak
Current: \url{https://github.com/alectionary/vibe-math/blob/master/formalization/distance-preservation/FamilyTranspose.lean}\\
Immutable: \url{https://github.com/alectionary/vibe-math/blob/98eb91cd9049e6a88887a5ba79150ad96a4a3e53/formalization/distance-preservation/FamilyTranspose.lean}\par\end{minipage}\par
\par\smallskip\noindent\begin{minipage}{\linewidth}\small\setlength{\baselineskip}{11pt}
\textbf{Source file:} \nolinkurl{FamilyTreeStopping.lean}\par\nopagebreak
Current: \url{https://github.com/alectionary/vibe-math/blob/master/formalization/distance-preservation/FamilyTreeStopping.lean}\\
Immutable: \url{https://github.com/alectionary/vibe-math/blob/98eb91cd9049e6a88887a5ba79150ad96a4a3e53/formalization/distance-preservation/FamilyTreeStopping.lean}\par\end{minipage}\par
\par\smallskip\noindent\begin{minipage}{\linewidth}\small\setlength{\baselineskip}{11pt}
\textbf{Source file:} \nolinkurl{FamilyWords.lean}\par\nopagebreak
Current: \url{https://github.com/alectionary/vibe-math/blob/master/formalization/distance-preservation/FamilyWords.lean}\\
Immutable: \url{https://github.com/alectionary/vibe-math/blob/98eb91cd9049e6a88887a5ba79150ad96a4a3e53/formalization/distance-preservation/FamilyWords.lean}\par\end{minipage}\par
\par\smallskip\noindent\begin{minipage}{\linewidth}\small\setlength{\baselineskip}{11pt}
\textbf{Source file:} \nolinkurl{FinalAudit.lean}\par\nopagebreak
Current: \url{https://github.com/alectionary/vibe-math/blob/master/formalization/distance-preservation/FinalAudit.lean}\\
Immutable: \url{https://github.com/alectionary/vibe-math/blob/98eb91cd9049e6a88887a5ba79150ad96a4a3e53/formalization/distance-preservation/FinalAudit.lean}\par\end{minipage}\par
\par\smallskip\noindent\begin{minipage}{\linewidth}\small\setlength{\baselineskip}{11pt}
\textbf{Source file:} \nolinkurl{GeneralLP.lean}\par\nopagebreak
Current: \url{https://github.com/alectionary/vibe-math/blob/master/formalization/distance-preservation/GeneralLP.lean}\\
Immutable: \url{https://github.com/alectionary/vibe-math/blob/98eb91cd9049e6a88887a5ba79150ad96a4a3e53/formalization/distance-preservation/GeneralLP.lean}\par\end{minipage}\par
\par\smallskip\noindent\begin{minipage}{\linewidth}\small\setlength{\baselineskip}{11pt}
\textbf{Source file:} \nolinkurl{GeneralLPPhysical.lean}\par\nopagebreak
Current: \url{https://github.com/alectionary/vibe-math/blob/master/formalization/distance-preservation/GeneralLPPhysical.lean}\\
Immutable: \url{https://github.com/alectionary/vibe-math/blob/98eb91cd9049e6a88887a5ba79150ad96a4a3e53/formalization/distance-preservation/GeneralLPPhysical.lean}\par\end{minipage}\par
\par\smallskip\noindent\begin{minipage}{\linewidth}\small\setlength{\baselineskip}{11pt}
\textbf{Source file:} \nolinkurl{Hypotheses11.lean}\par\nopagebreak
Current: \url{https://github.com/alectionary/vibe-math/blob/master/formalization/distance-preservation/Hypotheses11.lean}\\
Immutable: \url{https://github.com/alectionary/vibe-math/blob/98eb91cd9049e6a88887a5ba79150ad96a4a3e53/formalization/distance-preservation/Hypotheses11.lean}\par\end{minipage}\par
\par\smallskip\noindent\begin{minipage}{\linewidth}\small\setlength{\baselineskip}{11pt}
\textbf{Source file:} \nolinkurl{ParametricProjection.lean}\par\nopagebreak
Current: \url{https://github.com/alectionary/vibe-math/blob/master/formalization/distance-preservation/ParametricProjection.lean}\\
Immutable: \url{https://github.com/alectionary/vibe-math/blob/98eb91cd9049e6a88887a5ba79150ad96a4a3e53/formalization/distance-preservation/ParametricProjection.lean}\par\end{minipage}\par
\par\smallskip\noindent\begin{minipage}{\linewidth}\small\setlength{\baselineskip}{11pt}
\textbf{Source file:} \nolinkurl{ParametricRing.lean}\par\nopagebreak
Current: \url{https://github.com/alectionary/vibe-math/blob/master/formalization/distance-preservation/ParametricRing.lean}\\
Immutable: \url{https://github.com/alectionary/vibe-math/blob/98eb91cd9049e6a88887a5ba79150ad96a4a3e53/formalization/distance-preservation/ParametricRing.lean}\par\end{minipage}\par
\par\smallskip\noindent\begin{minipage}{\linewidth}\small\setlength{\baselineskip}{11pt}
\textbf{Source file:} \nolinkurl{ParametricSource.lean}\par\nopagebreak
Current: \url{https://github.com/alectionary/vibe-math/blob/master/formalization/distance-preservation/ParametricSource.lean}\\
Immutable: \url{https://github.com/alectionary/vibe-math/blob/98eb91cd9049e6a88887a5ba79150ad96a4a3e53/formalization/distance-preservation/ParametricSource.lean}\par\end{minipage}\par
\par\smallskip\noindent\begin{minipage}{\linewidth}\small\setlength{\baselineskip}{11pt}
\textbf{Source file:} \nolinkurl{ParametricSpecialization.lean}\par\nopagebreak
Current: \url{https://github.com/alectionary/vibe-math/blob/master/formalization/distance-preservation/ParametricSpecialization.lean}\\
Immutable: \url{https://github.com/alectionary/vibe-math/blob/98eb91cd9049e6a88887a5ba79150ad96a4a3e53/formalization/distance-preservation/ParametricSpecialization.lean}\par\end{minipage}\par
\par\smallskip\noindent\begin{minipage}{\linewidth}\small\setlength{\baselineskip}{11pt}
\textbf{Source file:} \nolinkurl{ParametricWords.lean}\par\nopagebreak
Current: \url{https://github.com/alectionary/vibe-math/blob/master/formalization/distance-preservation/ParametricWords.lean}\\
Immutable: \url{https://github.com/alectionary/vibe-math/blob/98eb91cd9049e6a88887a5ba79150ad96a4a3e53/formalization/distance-preservation/ParametricWords.lean}\par\end{minipage}\par
\par\smallskip\noindent\begin{minipage}{\linewidth}\small\setlength{\baselineskip}{11pt}
\textbf{Source file:} \nolinkurl{QuantumRows11.lean}\par\nopagebreak
Current: \url{https://github.com/alectionary/vibe-math/blob/master/formalization/distance-preservation/QuantumRows11.lean}\\
Immutable: \url{https://github.com/alectionary/vibe-math/blob/98eb91cd9049e6a88887a5ba79150ad96a4a3e53/formalization/distance-preservation/QuantumRows11.lean}\par\end{minipage}\par
\par\smallskip\noindent\begin{minipage}{\linewidth}\small\setlength{\baselineskip}{11pt}
\textbf{Source file:} \nolinkurl{QuantumWitness11.lean}\par\nopagebreak
Current: \url{https://github.com/alectionary/vibe-math/blob/master/formalization/distance-preservation/QuantumWitness11.lean}\\
Immutable: \url{https://github.com/alectionary/vibe-math/blob/98eb91cd9049e6a88887a5ba79150ad96a4a3e53/formalization/distance-preservation/QuantumWitness11.lean}\par\end{minipage}\par
\par\smallskip\noindent\begin{minipage}{\linewidth}\small\setlength{\baselineskip}{11pt}
\textbf{Source file:} \nolinkurl{SmallStopping.lean}\par\nopagebreak
Current: \url{https://github.com/alectionary/vibe-math/blob/master/formalization/distance-preservation/SmallStopping.lean}\\
Immutable: \url{https://github.com/alectionary/vibe-math/blob/98eb91cd9049e6a88887a5ba79150ad96a4a3e53/formalization/distance-preservation/SmallStopping.lean}\par\end{minipage}\par
\par\smallskip\noindent\begin{minipage}{\linewidth}\small\setlength{\baselineskip}{11pt}
\textbf{Source file:} \nolinkurl{Source11.lean}\par\nopagebreak
Current: \url{https://github.com/alectionary/vibe-math/blob/master/formalization/distance-preservation/Source11.lean}\\
Immutable: \url{https://github.com/alectionary/vibe-math/blob/98eb91cd9049e6a88887a5ba79150ad96a4a3e53/formalization/distance-preservation/Source11.lean}\par\end{minipage}\par
\par\smallskip\noindent\begin{minipage}{\linewidth}\small\setlength{\baselineskip}{11pt}
\textbf{Source file:} \nolinkurl{SourceCoordinates.lean}\par\nopagebreak
Current: \url{https://github.com/alectionary/vibe-math/blob/master/formalization/distance-preservation/SourceCoordinates.lean}\\
Immutable: \url{https://github.com/alectionary/vibe-math/blob/98eb91cd9049e6a88887a5ba79150ad96a4a3e53/formalization/distance-preservation/SourceCoordinates.lean}\par\end{minipage}\par
\par\smallskip\noindent\begin{minipage}{\linewidth}\small\setlength{\baselineskip}{11pt}
\textbf{Source file:} \nolinkurl{SourceDistance.lean}\par\nopagebreak
Current: \url{https://github.com/alectionary/vibe-math/blob/master/formalization/distance-preservation/SourceDistance.lean}\\
Immutable: \url{https://github.com/alectionary/vibe-math/blob/98eb91cd9049e6a88887a5ba79150ad96a4a3e53/formalization/distance-preservation/SourceDistance.lean}\par\end{minipage}\par
\par\smallskip\noindent\begin{minipage}{\linewidth}\small\setlength{\baselineskip}{11pt}
\textbf{Source file:} \nolinkurl{SourceModel.lean}\par\nopagebreak
Current: \url{https://github.com/alectionary/vibe-math/blob/master/formalization/distance-preservation/SourceModel.lean}\\
Immutable: \url{https://github.com/alectionary/vibe-math/blob/98eb91cd9049e6a88887a5ba79150ad96a4a3e53/formalization/distance-preservation/SourceModel.lean}\par\end{minipage}\par
\par\smallskip\noindent\begin{minipage}{\linewidth}\small\setlength{\baselineskip}{11pt}
\textbf{Source file:} \nolinkurl{SourceSectors.lean}\par\nopagebreak
Current: \url{https://github.com/alectionary/vibe-math/blob/master/formalization/distance-preservation/SourceSectors.lean}\\
Immutable: \url{https://github.com/alectionary/vibe-math/blob/98eb91cd9049e6a88887a5ba79150ad96a4a3e53/formalization/distance-preservation/SourceSectors.lean}\par\end{minipage}\par
\par\smallskip\noindent\begin{minipage}{\linewidth}\small\setlength{\baselineskip}{11pt}
\textbf{Source file:} \nolinkurl{SourceSymmetry.lean}\par\nopagebreak
Current: \url{https://github.com/alectionary/vibe-math/blob/master/formalization/distance-preservation/SourceSymmetry.lean}\\
Immutable: \url{https://github.com/alectionary/vibe-math/blob/98eb91cd9049e6a88887a5ba79150ad96a4a3e53/formalization/distance-preservation/SourceSymmetry.lean}\par\end{minipage}\par
\par\smallskip\noindent\begin{minipage}{\linewidth}\small\setlength{\baselineskip}{11pt}
\textbf{Source file:} \nolinkurl{SourceWordBridge.lean}\par\nopagebreak
Current: \url{https://github.com/alectionary/vibe-math/blob/master/formalization/distance-preservation/SourceWordBridge.lean}\\
Immutable: \url{https://github.com/alectionary/vibe-math/blob/98eb91cd9049e6a88887a5ba79150ad96a4a3e53/formalization/distance-preservation/SourceWordBridge.lean}\par\end{minipage}\par
\par\smallskip\noindent\begin{minipage}{\linewidth}\small\setlength{\baselineskip}{11pt}
\textbf{Source file:} \nolinkurl{TannerSix.lean}\par\nopagebreak
Current: \url{https://github.com/alectionary/vibe-math/blob/master/formalization/distance-preservation/TannerSix.lean}\\
Immutable: \url{https://github.com/alectionary/vibe-math/blob/98eb91cd9049e6a88887a5ba79150ad96a4a3e53/formalization/distance-preservation/TannerSix.lean}\par\end{minipage}\par
\par\smallskip\noindent\begin{minipage}{\linewidth}\small\setlength{\baselineskip}{11pt}
\textbf{Source file:} \nolinkurl{TransposeConstant.lean}\par\nopagebreak
Current: \url{https://github.com/alectionary/vibe-math/blob/master/formalization/distance-preservation/TransposeConstant.lean}\\
Immutable: \url{https://github.com/alectionary/vibe-math/blob/98eb91cd9049e6a88887a5ba79150ad96a4a3e53/formalization/distance-preservation/TransposeConstant.lean}\par\end{minipage}\par
\par\smallskip\noindent\begin{minipage}{\linewidth}\small\setlength{\baselineskip}{11pt}
\textbf{Source file:} \nolinkurl{TransposeDistance.lean}\par\nopagebreak
Current: \url{https://github.com/alectionary/vibe-math/blob/master/formalization/distance-preservation/TransposeDistance.lean}\\
Immutable: \url{https://github.com/alectionary/vibe-math/blob/98eb91cd9049e6a88887a5ba79150ad96a4a3e53/formalization/distance-preservation/TransposeDistance.lean}\par\end{minipage}\par
\par\smallskip\noindent\begin{minipage}{\linewidth}\small\setlength{\baselineskip}{11pt}
\textbf{Source file:} \nolinkurl{VoltageCycles.lean}\par\nopagebreak
Current: \url{https://github.com/alectionary/vibe-math/blob/master/formalization/distance-preservation/VoltageCycles.lean}\\
Immutable: \url{https://github.com/alectionary/vibe-math/blob/98eb91cd9049e6a88887a5ba79150ad96a4a3e53/formalization/distance-preservation/VoltageCycles.lean}\par\end{minipage}\par

